\chapter{Eco-Genoma Primordial: La Firma de la Creación}
El universo no solo posee una historia, sino una memoria inscrita en su estructura más íntima. A esta huella digital de las condiciones iniciales le denominamos \emph{Eco-Genoma Primordial}.

\section{Analogía con el Genoma Biológico}
Si el genoma biológico es un libro de instrucciones para la vida, el Eco-Genoma es el libro de instrucciones para la realidad misma. Mientras el genoma biológico está sujeto a mutaciones, el Eco-Genoma Primordial es \emph{estructuralmente invariante}.

\section{Tabla Periódica Informacional (TPI)}
La TPI clasifica los \emph{elementos de información} en tres dimensiones fundamentales:
\begin{itemize}
    \item \textbf{Dimensión de Coherencia (C):} Grado de organización interna.
    \item \textbf{Dimensión de Complejidad (K):} Cantidad de información necesaria para describir el patrón.
    \item \textbf{Dimensión de Resonancia (R):} Capacidad para interactuar y sincronizarse.
\end{itemize}

\section{Estructura del Reporte e01}
\begin{table}[H]
\centering
\begin{tabular}{lll}
\toprule
\textbf{Parámetro} & \textbf{Valor} & \textbf{Descripción} \\
\midrule
ID de Eco & e01 & Identificador del linaje. \\
Linaje & Patrón Grafeno & Estructura reticular hexagonal. \\
Resonancia & 0.855 & Sincronización global. \\
Fidelidad & 100\% & Reproducibilidad del patrón. \\
\bottomrule
\end{tabular}
\caption{Huella Genética de la Creación (Reporte e01)}
\end{table}

\section{El Patrón Grafeno Primordial}
La configuración inicial de máxima coherencia resultó ser una red reticular hexagonal en 6D. Esto sugiere que el universo, en su nivel más fundamental, prefiere la geometría del panal de abejas para tejer su existencia.
